% ECONOMICS_PROBLEM: {"id": "F12-161", "title": "Full gains from trade with heterogeneity, innovation and misallocation", "jel_code": "F12", "source_id": "OP-0081"}
% !TeX program = xelatex
\documentclass[11pt,a4paper]{article}
\usepackage[margin=23mm]{geometry}
\usepackage{fontspec}
\setmainfont{Latin Modern Roman}
\usepackage{amsmath,amssymb,mathtools}
\usepackage{xcolor,enumitem,booktabs,longtable,array,xurl,hyperref}
\definecolor{muted}{HTML}{53636C}
\hypersetup{colorlinks=true,urlcolor=blue,linkcolor=blue}
\setlength{\parindent}{0pt}
\setlength{\parskip}{5pt}
\newcommand{\R}{\mathbb R}
\newcommand{\E}{\mathbb E}
\newcommand{\N}{\mathbb N}
\newcommand{\ind}{\mathbf 1}
\newcommand{\Var}{\operatorname{Var}}
\newcommand{\Cov}{\operatorname{Cov}}
\newcommand{\supp}{\operatorname{supp}}
\DeclareMathOperator*{\argmin}{arg\,min}
\DeclareMathOperator*{\argmax}{arg\,max}
\newcommand{\entrymeta}[3]{{\sffamily\small\color{muted}\textbf{\texttt{#1}}\quad|\quad #2\par #3\par}}
\newcommand{\Problemref}[1]{\texttt{#1}}
\newcommand{\JELref}[2]{\texttt{#2} (JEL \texttt{#1})}
\begin{document}
\section*{Full gains from trade with heterogeneity, innovation and misallocation}
\textbf{Problem ID:} \texttt{F12-161}\quad \textbf{JEL:} \texttt{F12}\par
% BEGIN ECONOMICS STATEMENT
\label{problem:OP-0081}
\label{audited:ECON-079}
\label{atlas:prob:T1}

\noindent\textbf{Atlas ID:} T1; \textbf{original number:} 81; \textbf{field:} International trade. \textbf{Classification:} Editorial research formulation (theoretical/empirical); not a certified unsolved theorem.

\subsection*{Definitions and assumptions}
Fix a finite set of countries and an infinite sequence of dates. A structural model $\theta$ specifies household preferences, endowments, heterogeneous firms, production sets, innovation and entry costs, and domestic wedges. Firms choose feasible inputs, export destinations, entry and innovation; households choose budget-feasible consumption, labor and saving. Trade policy $\tau$ specifies bilateral barriers and tariff-revenue redistribution. Let $e_\theta(\tau)$ be a specified competitive or monopolistically competitive equilibrium selection with market clearing and consistent state transitions. For household $h$, let $c_{ht}^\tau>0$ and $\ell_{ht}^\tau$ be equilibrium consumption and labor, with utility $u_h(c,\ell)$ strictly increasing in consumption and assume $\mathbb E_\theta\sum_{t\ge0}\beta_h^t|u_h(c_{ht}^\tau,\ell_{ht}^\tau)|<\infty$, where $\beta_h\in(0,1)$.

\subsection*{Formal research question}
For policies $\tau^0,\tau^1$, define the consumption-equivalent gain $g_h(\theta)$ by
\[
 \mathbb E_\theta\sum_{t\ge0}\beta_h^tu_h((1+g_h)c_{ht}^{\tau^0},\ell_{ht}^{\tau^0})
 =\mathbb E_\theta\sum_{t\ge0}\beta_h^tu_h(c_{ht}^{\tau^1},\ell_{ht}^{\tau^1}).
\]
Require $g_h>-1$ and strict monotonicity and range conditions ensuring existence and uniqueness. Labor and its disutility are held at the respective equilibrium paths in this comparison. Let $\Theta(P_O)$ be a prespecified class of models and selections generating the observed law $P_O$ of trade, firm, productivity and innovation data. Characterize the identified set of household gains and an explicitly weighted welfare aggregate. Determine which additional moments or interventions shrink it when innovation and domestic wedges respond endogenously to trade.

\subsection*{Interpretation and scope}
The objective is model comparison and identification, not an assumption that all matching models yield the same gain. Policy comparisons must hold the initial state and exogenous technology shocks fixed while permitting endogenous innovation and reallocation. A meaningful result may provide welfare bounds, testable restrictions distinguishing mechanisms, or a documented quantitative estimate within a stated model. Distributional gains cannot be inferred from aggregate trade shares alone. Report sensitivity to transition costs, revenue recycling, financing constraints and the equilibrium selection. Validation should include firm and innovation moments excluded from estimation.

\subsection*{Source audit and status}
[S1] and [S2] establish benchmark trade mechanisms. [S3] establishes a sufficient-statistic relationship under its maintained restrictions; it is not a theorem for arbitrary endogenous innovation or misallocation. This entry is an editorial extension of those benchmarks. Their publication does not certify the extension as an unsolved mathematical problem in 2026.

\subsection*{References}

\begin{enumerate}[label={[S\arabic*]},leftmargin=*]

\item Jonathan Eaton and Samuel Kortum (2002). \emph{Technology, Geography, and Trade}. Econometrica 70(5), 1741–1779. \url{https://doi.org/10.1111/1468-0262.00352}.
\textit{Locator and role:} Publisher or author abstract / bibliographic record; Ricardian general equilibrium, geographic trade costs and counterfactual gains. Provides background or a benchmark result; does not establish that the editorial question remains unresolved in 2026.

\item Marc J. Melitz (2003). \emph{The Impact of Trade on Intra-Industry Reallocations and Aggregate Industry Productivity}. Econometrica 71(6), 1695–1725. \url{https://doi.org/10.1111/1468-0262.00467}.
\textit{Locator and role:} Publisher or author abstract / bibliographic record; Firm selection and trade-induced productivity reallocation. Provides background or a benchmark result; does not establish that the editorial question remains unresolved in 2026.

\item Costas Arkolakis, Arnaud Costinot, and Andrés Rodríguez-Clare (2012). \emph{New Trade Models, Same Old Gains?}. American Economic Review 102(1), 94–130. \url{https://doi.org/10.1257/aer.102.1.94}.
\textit{Locator and role:} Publisher or author abstract / bibliographic record; Sufficient-statistic welfare comparisons within the paper’s restricted model class. Provides background or a benchmark result; does not establish that the editorial question remains unresolved in 2026.

\end{enumerate}

\subsection*{Integrated refinement: ECON-079}
{\sffamily\small\color{muted}Trade gains under weak institutions and heterogeneous constraints\par}
This refinement adopts the additional source conventions
(page~\pageref{audited:conventions}), subject to its local restrictions.
\textbf{Institutional enforcement and heterogeneous financing constraints.}
Add a declared contract-enforcement state $\kappa_r$ and firm financing
constraints $b_i$ to the trade-gain model. For the same market-access regime,
characterize conditional welfare effects
$\theta(\kappa)=\mathbb E[W_r(1)-W_r(0)\mid\kappa_r=\kappa]$ and their compatible
sets. Define $W_r$ as real income or a specific welfare measure and include
consumer prices, firm profits and worker earnings without double counting.
Separate pretreatment institutional effect modifiers from policy-induced
institutions. Identifying a financing/enforcement mediation channel requires
credible channel interventions or restrictions beyond conditional correlations.
Report support and invariance assumptions before transporting effects across
economies with different institutional constraints.
\subsection*{Supplied source audit for ECON-079}
Local [S1], [S2], etc. in this audit and the references immediately below
belong to ECON-079, independently of the original entry's reference numbering.
[S1] and [S2] support trade with institutional frictions and heterogeneous constraints. This baseline-conditional welfare target is editorial; the review supplies no universal weak-institution treatment effect.
\subsection*{References for integrated refinement ECON-079}
{\footnotesize\raggedright
\par\noindent\textbf{[S1]} David Atkin, Amit Khandelwal, Laura Boudreau, Rafael Dix-Carneiro, Isabela Manelici, Pamela Medina, Brian McCaig, Ameet Morjaria, Luigi Pascali, Heitor Pellegrina, Bob Rijkers, and Meredith Startz (2025). \emph{International Trade}. VoxDevLit 4(2), February 2025.\par \textit{Verified locator / source role:} Primary publisher page checked for review identity, authors, version and background. The specific published agenda is corroborated in [S2]; the exact formal target is editorial.\par \url{https://voxdev.org/voxdevlit/international-trade}\par\smallskip
\par\noindent\textbf{[S2]} Oliver Hanney / VoxDev (living compilation). \emph{Evidence gaps: Key unanswered questions in development economics}. Accessed 8 October 2026. The page currently covers 20 topics; the ``16-topics URL is a historical slug.\par \textit{Verified locator / source role:} Topic heading: International Trade. Supports the broad research agenda; this is not a verbatim quotation or an attribution of the displayed mathematical target.\par \url{https://voxdev.org/topic/evidence-gaps-key-unanswered-questions-across-16-topics-development-economics}\par\smallskip
}

\subsection*{Common conventions: Source mathematical conventions}

Unless an entry states otherwise, agent and firm sets are finite. State and action spaces are measurable, strategies and outcome maps are measurable, probability laws are specified, and expectations exist for the targets under discussion. Infinite-horizon discount factors lie in $(0,1)$ and discounted objectives are finite. Norms, tolerances and complexity input sizes must be specified for computational comparisons. Constants or primitives that appear locally are inputs, not empirical facts asserted by this catalogue.

\textbf{Identification.} A statistical model $m\in\mathcal M$ implies an observable law $P_m$ and target $\theta(m)$. At law $P$, its identified set is
\[
\Theta(P;\mathcal M)=\{\theta(m):m\in\mathcal M,\ P_m=P\}.
\]
Point identification holds when this set is a singleton, or equivalently when $P_{m_1}=P_{m_2}$ implies $\theta(m_1)=\theta(m_2)$ for all compatible models. Sharp bounds mean bounds attained, or approached when the boundary is not attained, by compatible models. Writing this set is a definition, not a solution: the substantive task is to establish informative restrictions and derive or estimate the set. An unrestricted-confounding model may give only trivial bounds.

\textbf{Inference.} Identification, estimability and valid uncertainty quantification are separate questions. Fix $\alpha\in(0,1)$. A confidence procedure $C_n$ over a declared law class $\mathcal P$ has asymptotic uniform coverage at level $1-\alpha$ when
\[
\liminf_{n\to\infty}\inf_{P\in\mathcal P}\Pr_P\{\theta(P)\in C_n\}\geq1-\alpha.
\]
For set-identified targets the coverage event must be defined separately, for example coverage of the full identified set. If an entry uses identification-indexed models instead of uniquely indexed laws, the same coverage convention is applied over those models. Pointwise limits do not establish uniform validity, and asymptotic validity does not guarantee finite-sample precision. Sample sizes, dependence structures and weak-identification sequences are part of each application.

\textbf{Counterfactuals.} $Y_i(a)$ is a potential outcome under a specified intervention, including its timing, implementation and versions. It is not $Y_i$ conditional on voluntarily choosing $a$. A full intervention vector permits interference; shorthand scalar treatment notation does not silently impose no interference. Assignment, selection, attrition, anticipation and measurement must be specified. A claim about an instrument's local effect is not automatically a population policy effect.

\textbf{Economic equilibrium.} A counterfactual model includes best responses, constraints, feasibility, market clearing, state transitions and expectations consistent with the stated information. When several equilibria exist, either a declared selector is an input or the target is set-valued. Comparisons across policy regimes must use the stated selection convention. Reduced-form relationships are not presumed invariant after an equilibrium-changing intervention.

\textbf{Optimization and welfare.} Optimization is over the stated feasible set. If attainment is not established, $\sup$ or $\inf$ and $\varepsilon$-optimal choices replace an asserted maximizer or minimizer. For example, with value $V=\sup_{a\in\mathcal A}W(a)<\infty$, an $\varepsilon$-optimal set is $\{a\in\mathcal A:W(a)\geq V-\varepsilon\}$ for $\varepsilon>0$. Benefits, costs, risk and health or environmental outcomes can be aggregated only using declared units and conversion weights. Interpersonal utility weights, the treatment of fiscal transfers and foreign profits, discounting, and the behavioral welfare criterion are explicit inputs. A comparative-static derivative additionally requires existence and appropriate differentiability of the selected counterfactual.

\textbf{Sources.} Local labels [S1], [S2], and so forth restart in every question. Locators identify the relevant abstract, model, section or result at the level actually checked; a bibliographic or abstract-level verification is not represented as inspection of an entire proof. Working papers are distinguished from later publications and changing versions. Source summaries are paraphrased. All URL checks and editorial formulations refer to the audit date stated above.

\subsection*{Common conventions: Additional-source status and mathematical conventions}

\label{audited:conventions}

Of the 100 supplied ECON entries, 65 distinct problems appear here; 32 close
matches refine earlier entries, and three duplicate questions map to existing
entries. Appendix~\ref{app:audited-crosswalk} lists every destination.
The attachment's P/R/S/U distinctions and source audits are retained. Its
precise mathematical formalizations are editorial unless the individual audit
says otherwise. Candidate subsidy targets ECON-004--006 retain their U flags;
the resolved approval-core question ECON-007 updates OP-0202 above.
The following supplied conventions govern both this part and the attributed
ECON refinements elsewhere in the catalogue, subject to local restrictions.

\textbf{Parameterized questions.} A problem family lists inputs that must be fixed before an instance is evaluated: population, horizon, policies, information, data law, model class and welfare weights. These are required choices, not quantities the citations are assumed to specify. Undefined applications should not be treated as identified estimands.

\textbf{Indices and integrability.} Sets of agents, firms and regions are finite unless a population measure is explicitly used. \(i,j\) index units, \(r\) regions, \(t\) dates and \(h\) horizons. \(\mathbb E\) and \(\Pr\) refer to the declared population and law; \(\mathbf1\{A\}\) is an indicator. Every displayed expectation and discounted sum needs existence in the stated domain. Infinite horizons use a discount in \((0,1)\) and a convergence condition; finite horizons may use discount one.

\textbf{Optimization and existence.} A displayed \(\max\), \(\min\), or \(\arg\max\) is conditional on attainment. For finite feasible menus this is immediate; general spaces need continuity and compactness/coercivity or another existence argument. Without attainment, the target is the supremum/infimum and arbitrarily near-optimal feasible choices. \(\inf\varnothing=+\infty\). Empty feasibility and an unidentified target are different issues.

\textbf{Potential outcomes.} \(Y_i(z)\) is the outcome under a specified intervention, not the observed conditional mean among people choosing \(z\). In binary comparisons \(\Delta Y_i=Y_i(1)-Y_i(0)\). Specify assignment, treatment versions, compliance, information and observation horizon. Interference is allowed under a full policy vector; compressed exposure notation needs a justified mapping. No interference, consistency and overlap are substantive assumptions, not automatic consequences of notation.

\textbf{Populations and observation.} Fix baseline eligibility and population weights. Conditioning on post-policy employment, survival, relocation or program uptake can change the population being compared. Death is not ordinary loss to follow-up; outcomes truncated by death need a composite convention, joint survival target, or explicitly defined survivor stratum. Fertility-changing interventions need a population functional rather than an unexplained fixed descendant cohort. Measurement and missingness models must be supplied.

\textbf{Identification.} A structural model \(m\in\mathcal M\) implies observable law \(P_m\) and target \(\theta(m)\). Identification means
\[
 P_{m_1}=P_{m_2}\ \Longrightarrow\ \theta(m_1)=\theta(m_2)
 \quad\text{for all }m_1,m_2\in\mathcal M .
\]
Without identification, the target is
\[
 \Theta(P;\mathcal M)=\{\theta(m):m\in\mathcal M,\ P_m=P\}.
\]
Writing \(\theta(P)\) before establishing identification can hide incompatible latent models. Confidence sets additionally need a sampling law, confidence level, asymptotic sequence and uniformity class. Point identification, coverage and informative precision are separate properties.

\textbf{Mediation.} Total effects of randomized assignment are distinct from cross-world natural indirect effects. Belief/effort-channel effects require a meaningful mediator intervention and additional measurement, exclusion or mediation assumptions. Randomization of the initial policy alone does not supply those assumptions. Report bounds or interventional alternatives when the stronger target is unsupported.

\textbf{Equilibrium.} A counterfactual \(W(z)\), \(p(z)\), or \(\mu_t^z\) requires individual optimization, feasibility, government/resource budgets, market clearing and state-transition laws. State equilibrium existence and selection, or report ranges over compatible equilibria. Initial-state integration includes subsequent endogenous transitions; current-state integration must use the policy-dependent distribution. Reduced-form invariance after policy is not assumed.

\textbf{Welfare accounts.} Scores, utility, health, dollars and ecological quantities require explicit conversion before addition. Fees, taxes, subsidies, wages and extortion payments are transfers between parties; distributional weights can make them welfare relevant, but their face value is not automatically a resource loss. State ownership, geographic boundaries, foreign-income weights and fiscal finance. Do not add profits to household welfare already including dividends, or count land capitalization and the underlying benefit twice. A benefit-cost expression must distinguish gross benefits from net welfare and subtract every real cost exactly once.

\textbf{Derivatives and risk.} Log derivatives require positive arguments; local responses require admissible perturbations and specified equilibrium branches. Differentiation under expectations needs regularity/domination, and boundaries may allow only one-sided derivatives. Risk uses a specified law; ambiguity uses a declared nonempty law set on a common history space. Adaptive policies must be nonanticipative. A robust criterion is an editorial choice unless attributed explicitly.

\textbf{Scope overlaps.} Allocation, collective choice, contracts, household bargaining and coordinated policy overlap economics with optimization and game theory. They are retained because they appeared in the original catalogue; no claim of a strictly game-theory-free selection is made.
% END ECONOMICS STATEMENT
\end{document}
